\documentclass[11pt]{amsart}

\usepackage[margin=1in]{geometry}
\usepackage{amsmath,amssymb,amsthm}
\usepackage{booktabs}
\usepackage{enumitem}
\usepackage{hyperref}

\hypersetup{
  colorlinks=true,
  linkcolor=blue,
  citecolor=blue,
  urlcolor=blue
}

\newtheorem{theorem}{Theorem}[section]
\newtheorem{lemma}[theorem]{Lemma}
\newtheorem{proposition}[theorem]{Proposition}
\newtheorem{corollary}[theorem]{Corollary}
\theoremstyle{definition}
\newtheorem{definition}[theorem]{Definition}
\theoremstyle{remark}
\newtheorem{remark}[theorem]{Remark}

\newcommand{\F}{\mathbb{F}}
\newcommand{\Z}{\mathbb{Z}}
\newcommand{\Cay}{\operatorname{Cay}}
\newcommand{\Haar}{\mathcal{H}}
\newcommand{\ord}{\operatorname{ord}}
\newcommand{\cl}{\operatorname{cl}}

\title[Power-of-Two Cycles in Nilpotent Cubic Haar Graphs]{Power-of-Two Cycles in Cubic Haar Graphs over Nilpotent Groups of Class at Most Six}
\author{OvermanOu}
\email{ad110158@gmail.com}
\date{November 17, 2026}

\begin{document}

\begin{abstract}
The Erd\H{o}s--Gy\'arf\'as conjecture asserts that every finite simple graph of minimum degree at least three contains a simple cycle whose length is a power of two. Recently, Burkhart proved that every odd-order Cayley graph, and consequently every Cayley graph of a finite nilpotent group, satisfies the conjecture. However, cubic three-shift Haar graphs (bipartite bi-Cayley graphs) over odd-order or nonabelian nilpotent groups are generally non-Cayley, and their $3$-regularity rules out the single-vertex figure-eight induction central to Burkhart's argument.

Using exact integral Magnus expansions and Hall commutator collection in free rank-two nilpotent groups up to class six, we prove that every connected component of every simple cubic three-shift Haar graph over any finite group admitting a normal $2$-subgroup with nilpotent class-at-most-four quotient (unconditionally across all primes and exponents) contains a simple $2^k$-cycle ($k \ge 2$), tightening cleanly to $C_4$, $C_8$, $C_{16}$, or $C_{32}$ for all odd-order nilpotent groups of class at most four via $w_4$ ($p \ge 5$) and a two-stage commutator descent with algebraic rigidity at $p = 3$ down to the $9$-member critical family $G_{\mathrm{crit}}(\alpha,\beta)$. Extending the Magnus calculus to degree six shows that all three determinant-$27$ $64$-color words $w_{64}^{(1)}, w_{64}^{(2)}, w_{64}^{(3)}$ vanish identically modulo $\gamma_7(F_2)$, unconditionally resolving nilpotency classes $5$ and $6$ whenever the quotient's Sylow $3$-subgroup has class at most four (in particular, all finite $p$-groups of class at most $6$ with $p \ge 5$ and arbitrary exponent, verified on $UT_7(\F_5)$ of order $5^{21}$, class $6$, and exponent $25$). Conversely, exhaustive Magnus evaluation of all $288$ free-class-$4$ $32$-color words shows $0/288$ close at degree $5$, proving $64$ is minimal for free class $\ge 5$. Furthermore, $w_{64}^{(1)}$ exhibits all $992$ integral degree-$6$ prefix-difference gcds equal to $1$, yielding a single universal $C_{64}$ certificate across $(\Z/p^k\Z)\langle X,Y\rangle/(X,Y)^7$ for every prime $p \ge 2$ and $k \ge 1$, while $UT_4(\F_p)$ ($p \ge 7$, girth $14$) and $UT_6(\F_p)/\langle I+E_{16}\rangle$ ($p \ge 11$, girth $\ge 18$) demonstrate sharpness of $C_{32}$ in class $\le 4$. These results yield cover obstructions for connected graphs containing $\Theta(2^k,2^k,2^k)$.
\end{abstract}

\maketitle

\section{Introduction}

In 1995, Erd\H{o}s and Gy\'arf\'as~\cite{Erdos1997} conjectured that every finite simple graph $G$ with minimum degree $\delta(G) \ge 3$ contains a simple cycle of length $2^k$ for some integer $k \ge 2$. While Liu and Montgomery~\cite{LiuMontgomery2023} settled the conjecture for graphs of sufficiently large average degree, the small-degree regime---particularly $3$-regular (cubic) graphs---remains open. Recent computational and structural progress includes verified lower bounds of $24$ vertices for $\{4,8\}$-free graphs~\cite{Garcia2026}, $41$ vertices for general minimal counterexamples and $50$ for cubic graphs~\cite{DucoffeDumitru2026}, $60$ vertices for cubic bipartite graphs~\cite{Tranquilli2026}, and predominantly cubic structure theorems~\cite{Carr2026}.

On the algebraic side, Burkhart~\cite{Burkhart2026} recently proved that every Cayley graph of odd order with symmetric generating set $|S| \ge 3$ satisfies the Erd\H{o}s--Gy\'arf\'as conjecture, and deduced the same for Cayley graphs of all finite nilpotent groups. It is natural to ask how far algebraic cycle-forcing extends to the next most symmetric class of $3$-regular graphs: \emph{cubic Haar graphs} (also called $3$-regular bi-Cayley graphs).

\begin{definition}[Three-shift cubic Haar graph]
Let $\Gamma$ be a group and $S = \{s_0, s_1, s_2\} \subseteq \Gamma$ be a set of three distinct elements. The \emph{cubic Haar graph} $\Haar(\Gamma; S)$ is the simple $3$-regular bipartite graph with vertex set $\Gamma \times \{0,1\}$ and edge set
\[
E = \{ \{(g,0), (g s_c, 1)\} : g \in \Gamma,\ c \in \{0,1,2\} \}.
\]
\end{definition}

Cubic Haar graphs fall outside Burkhart's Cayley-graph theorem~\cite{Burkhart2026} for three structural reasons:
\begin{enumerate}[label=(\roman*)]
\item An odd-order group contains no involutions, so every undirected Cayley graph $\Cay(G, S)$ of odd order with $|S| \ge 3$ has even degree at least $4$; no cubic odd-order Cayley graph exists.
\item For a cubic Cayley graph over a finite nilpotent group $G$, the symmetric generating set $S$ of size $3$ must contain an involution, which necessarily lies in the normal Sylow $2$-subgroup $P_2 \trianglelefteq G$; consequently the odd quotient $G/P_2$ receives at most two generators, bypassing the nonabelian odd-order case in~\cite[Corollary~1.4]{Burkhart2026}. By contrast, $\Haar(K; \{1,a,b\})$ over an odd-order $p$-group $K$ is genuinely cubic, bipartite, and in general not a Cayley graph when $K$ is nonabelian.
\item Burkhart's inductive lifting lemma~\cite[Lemma~1.2]{Burkhart2026} lifts a \emph{symmetric figure-eight}---two power-of-two simple cycles intersecting at a single vertex. In any $3$-regular graph, two simple cycles meeting at a vertex $v$ must share at least one edge incident to $v$ (since $2 + 2 > 3$), so \emph{no figure-eight exists in any cubic graph}.
\end{enumerate}

A different mechanism is therefore needed to force power-of-two cycles in nonabelian cubic Haar graphs. We use exact integral noncommutative Magnus expansions and Hall commutator collection to obtain universal short closed color words whose same-shore prefix collisions are algebraically obstructed across entire nilpotent group families.

\subsection{Main results}

\begin{theorem}[Nilpotent class at most three]\label{thm:class3}
Let $\Gamma$ be a finite group having a normal $2$-subgroup $N$ such that $\Gamma/N$ is nilpotent of nilpotency class at most $3$ (in particular, $\Gamma$ itself may be any finite nilpotent group of class at most $3$). Then every connected component of every simple cubic three-shift Haar graph $\Haar(\Gamma; S)$ contains a simple cycle of length $2^k$ for some $k \ge 2$. Moreover, when $\Gamma$ is nilpotent of class at most $3$ and odd order, every connected component contains a simple cycle of length $4, 8, 16$, or $32$.
\end{theorem}

\begin{theorem}[Unconditional nilpotent class at most four and sharp odd-order $C_{32}$ bound]\label{thm:class4}
Let $\Gamma$ be any finite group admitting a normal $2$-subgroup $N \trianglelefteq \Gamma$ such that $\Gamma/N$ is nilpotent of nilpotency class at most $4$ (with no restriction on the Sylow $3$-subgroup or on exponents; in particular, $\Gamma$ itself may be any finite nilpotent group of class at most $4$). Then every connected component of every simple cubic three-shift Haar graph $\Haar(\Gamma; S)$ contains a simple cycle of length $2^k$ for some $k \ge 2$.

Moreover, whenever $\Gamma$ is a finite odd-order nilpotent group of nilpotency class at most $4$, every connected component of $\Haar(\Gamma; S)$ contains a simple cycle of length $4, 8, 16$, or $32$ (and $C_{32}$ is sharp against Proposition~\ref{prop:boundaries}).
\end{theorem}

Note that Theorem~\ref{thm:class4} requires neither nilpotency of $\Gamma$ itself nor any restriction on $\exp(\Gamma)$ or its $3$-Sylow class (covering nonnilpotent controls such as $P_7 \times A_4$ of order $28{,}812$, as well as $UT_5(\F_3)$ and $G_{\mathrm{crit}}(0,0) \times C_5$ of order $1{,}215$). Pushing the free-group Magnus calculus two steps higher yields:

\begin{theorem}[Nilpotency classes five and six]\label{thm:class6}
Let $\Gamma$ be any finite group admitting a normal $2$-subgroup $N \trianglelefteq \Gamma$ such that $Q = \Gamma/N$ is nilpotent of nilpotency class at most $6$ and its Sylow $3$-subgroup $P_3(Q)$ has nilpotency class at most $4$ (in particular, every finite $p$-group of class at most $6$ for any prime $p \ge 5$ and arbitrary exponent, and every mixed nilpotent group of class $\le 6$ with class-$\le 4$ $3$-Sylow). Then every connected component of $\Haar(\Gamma; S)$ contains a simple cycle of length $2^k$ ($k \ge 2$); moreover, when $N = 1$ (so $\Gamma$ has odd order), every component contains $C_4, C_8, C_{16}, C_{32}$, or $C_{64}$.
\end{theorem}

\begin{theorem}[Dyadic theta cover obstruction]\label{thm:theta}
Let $B$ be a connected finite simple graph containing a $\Theta(L,L,L)$ subgraph, i.e., three internally vertex-disjoint paths of common length $L = 2^k$ ($k \ge 1$) connecting two vertices of $B$. Let $\Gamma$ be any finite group satisfying the hypotheses of Theorem~\ref{thm:class3}, \ref{thm:class4}, or~\ref{thm:class6}. Then every connected component of every regular $\Gamma$-cover of $B$ contains a simple cycle whose length is a power of two.
\end{theorem}

Finally, we record strengthened all-prime $C_{64}$ certificates in degree-$6$ truncated free algebras together with the exact free-group length-$32$ barrier for class $\ge 5$ and sharpness controls:

\begin{proposition}[Degree-$6$ truncated free algebras and high-girth families]\label{prop:boundaries}
\leavevmode
\begin{enumerate}[label=(\alph*)]
\item For every prime $p \ge 2$ and every integer $k \ge 1$, let $A_{p,k}^{(6)} = (\Z/p^k\Z)\langle X,Y\rangle/(X,Y)^7$ and $K_{p,k}^{(6)} = \langle 1+X, 1+Y \rangle \le (A_{p,k}^{(6)})^\times$. Then the single word $w_{64}^{(1)}$ traces a simple $C_{64}$ in $\Haar(K_{p,k}^{(6)}; \{1, 1+X, 1+Y\})$ across all $p \ge 2$ and $k \ge 1$ (and hence also modulo $(X,Y)^5$).
\item For every prime $p \ge 7$, the two-generator subgroup $K = \langle I+E_{12}+E_{23}, I+E_{23}+E_{34}\rangle \le UT_4(\F_p)$ has order $p^4$ and class $3$, and $\Haar(K; \{I, A, B\})$ has girth $14$ (no $C_4, C_8, C_{16}$) and a simple $C_{32}$.
\item For every prime $p \ge 11$, the two-generator subgroup $K = \langle I+E_{12}+E_{34}+E_{56}, I+E_{23}+E_{45}\rangle \le UT_6(\F_p)/\langle I+E_{16}\rangle$ has order $p^6$ and class $4$, and $\Haar(K; \{I, A, B\})$ has girth at least $18$ (no $C_4, C_8, C_{16}$) and simple $C_{32}$ and $C_{64}$.
\end{enumerate}
\end{proposition}

\begin{proposition}[Free-group length-$32$ barrier for class at least five]\label{prop:barrier32}
Among all $288$ cyclically nonbacktracking $32$-color Haar words that close in $F_2/\gamma_5(F_2)$, exact integral evaluation in $\Z\langle X,Y\rangle/(X,Y)^6$ shows that $0/288$ close at degree $5$. Since $D_6(F_2) = \gamma_6(F_2)$, length $64$ is minimal for any universal class-$\ge 5$ free-group certificate.
\end{proposition}

\section{Preliminaries, Abelian and Class-$2$ Haar Graphs, and $2$-Kernel Lifting}\label{sec:prelim}

\subsection{Shift normalization and word conventions}

Given $\Haar(\Gamma; \{s_0, s_1, s_2\})$, the bijection $(g, 0) \mapsto (g s_0, 0)$, $(h, 1) \mapsto (h, 1)$ is a graph isomorphism onto $\Haar(\Gamma; \{1, a, b\})$ where $a = s_0^{-1} s_1$ and $b = s_0^{-1} s_2$ are distinct nonidentity elements. Let $K = \langle a, b \rangle \le \Gamma$. The connected component of $(1,0)$ is $\Haar(K; \{1, a, b\})$, and every connected component is a left coset translate $(gK \times \{0\}) \cup (gK \times \{1\}) \cong \Haar(K; \{1, a, b\})$. Thus it suffices to work inside $\Haar(K; \{1,a,b\})$ with $K = \langle a, b \rangle$.

Represent a closed walk starting at $(1,0)$ in $\Haar(K; \{1, a, b\})$ by an even-length cyclic word $w = c_0 c_1 \cdots c_{2m-1} \in \{0,1,2\}^{2m}$, where colors $0, 1, 2$ denote shifts $s_0 = 1, s_1 = a, s_2 = b$. At even step $i$ (shore $0$ to shore $1$) we right-multiply by $s_{c_i}$; at odd step $i$ (shore $1$ to shore $0$) we right-multiply by $s_{c_i}^{-1}$. Define prefix group coordinates $P_0 = 1$ and $P_{i+1} = P_i s_{c_i}^{(-1)^i}$ ($0 \le i < 2m$). The vertex at step $i$ is $(i \bmod 2, P_i)$. The walk is \emph{cyclically nonbacktracking} if $c_i \ne c_{i+1}$ for all $i \pmod{2m}$, closes if $P_{2m} = 1$, and forms a \emph{simple cycle} $C_{2m}$ if $P_{2m} = 1$ and $P_i^{-1} P_j \ne 1$ for all $0 \le i < j < 2m$ with $i \equiv j \pmod 2$.

\subsection{Abelian and class-two foundations}

\begin{theorem}[Universal abelian Haar obstruction]\label{thm:abelian}
Let $A$ be any abelian group (finite or infinite), and let $0, a, b \in A$ be three distinct elements. Then every connected component of $\Haar(A; \{0, a, b\})$ contains a simple $C_4$, $C_8$, or $C_{16}$.
\end{theorem}

\begin{proof}
By shift normalization we consider $S = \{0, a, b\}$ inside $\langle a, b \rangle \le A$, and examine the cyclically nonbacktracking $16$-color word
\[
w_{16} = \mathtt{0101012010201212}.
\]
In $\Z^2$ with $a=(1,0)$ and $b=(0,1)$, $w_{16}$ closes at $(0,0)$ and visits $8$ distinct vertices on each shore. Taking canonical differences $v \in \Z^2$ (up to sign, choosing $v > (0,0)$ lexicographically) across all $56$ same-shore prefix pairs gives exactly $13$ vectors: three collapse vectors $(0,1)$, $(1,-1)$, $(1,0)$ (excluded since $0,a,b$ are distinct), and ten non-collapse vectors realized as endpoints of the cyclically nonbacktracking short words in Table~\ref{tab:abelian_short}.

\begin{table}[ht]
\centering
\caption{Short cyclically nonbacktracking words covering all non-collapse collisions of $w_{16}$.}
\label{tab:abelian_short}
\begin{tabular}{cccc}
\toprule
Vector $v$ & Short word & Vector $v$ & Short word \\
\midrule
$(0,2)$  & \texttt{0202}     & $(2,-1)$ & \texttt{0121} \\
$(1,-2)$ & \texttt{0212}     & $(2,0)$  & \texttt{0101} \\
$(1,1)$  & \texttt{0102}     & $(2,1)$  & \texttt{01210202} \\
$(1,2)$  & \texttt{01010212} & $(3,-1)$ & \texttt{01212102} \\
$(2,-2)$ & \texttt{1212}     & $(3,0)$  & \texttt{01012102} \\
\bottomrule
\end{tabular}
\end{table}

If $w_{16}$ is simple in $A$, it yields $C_{16}$. Otherwise some $v$ from Table~\ref{tab:abelian_short} vanishes in $A$, so its short word of length $4$ or $8$ closes. In a simple bipartite graph, a $4$-step cyclically nonbacktracking closed walk is a simple $C_4$, and an $8$-step cyclically nonbacktracking closed walk is either a simple $C_8$ or splits at distance $4$ into a simple $C_4$.
\end{proof}

\begin{proposition}[Class-two Haar obstruction]\label{prop:class2}
Let $K = \langle a, b \rangle$ have nilpotency class at most $2$, and set $c = [a,b] = a^{-1} b^{-1} a b$. If $\ord(c) \ge 3$ (including $\infty$), then $\Haar(K; \{1, a, b\})$ contains a simple $C_{16}$ given by one of the two cyclically nonbacktracking words
\[
W_0 = \mathtt{0101201210102102},\qquad W_1 = \mathtt{0102012010121021}.
\]
\end{proposition}

\begin{proof}
In a class-$2$ group, $c$ is central and $a^x b^Y = b^Y a^x c^{xY}$, so the free class-$2$ group is parameterized by triples $(x,y,z) \in \Z^3$ representing $b^y a^x c^z$ with product $(x,y,z)(X,Y,Z) = (x+X, y+Y, z+Z+xY)$ and inverse $(-x,-y,-z+xy)$. Both $W_0$ and $W_1$ end at $(0,0,0)$ in the free class-$2$ group, hence close in $K$. Their $8$ prefix triples $(x,y,z)$ on each shore are pairwise distinct and satisfy $|\Delta x| \le 3$, $|\Delta y| \le 1$, and $|\Delta z| \le 2$ for every same-shore difference $(\Delta x, \Delta y, \Delta z) = (x_1 - x_2, y_1 - y_2, z_1 - z_2)$.

Suppose two same-shore prefixes $u=(x_1,y_1,z_1)$ and $v=(x_2,y_2,z_2)$ collide in $K$, so $v^{-1}u = 1$ in $K$. Since $v^{-1}u = b^{\Delta y} a^{\Delta x} c^{\Delta z'}$ in $K$, commuting with $a$ and $b$ yields $c^{-\Delta y} = [v^{-1}u, a] = 1$ and $c^{\Delta x} = [v^{-1}u, b] = 1$. Because $\ord(c) \ge 3$ and $|\Delta y| \le 1$, we must have $\Delta y = 0$; consequently $y_1 = y_2$ and $v^{-1} u = a^{\Delta x} c^{\Delta z}$ with $\Delta z = z_1 - z_2$.

If $\ord(c) \ge 4$ (or $\infty$), then $c^{\Delta x} = 1$ with $|\Delta x| \le 3$ forces $\Delta x = 0$, leaving $c^{\Delta z} = 1$ with $1 \le |\Delta z| \le 2$, a contradiction; thus $W_0$ is already a simple $C_{16}$. If $\ord(c) = 3$, collision forces $\Delta y = 0$ and $\Delta x \in \{0, 3, -3\}$. The case $\Delta x = 0$ again contradicts $1 \le |\Delta z| \le 2$, while $\Delta x = \pm 3$ requires $a^3 \in \langle c \rangle$, say $a^3 = c^h$ for a unique $h \in \{0,1,2\}$. Then $a^{\Delta x} c^{\Delta z} = c^{\Delta z + h(\Delta x/3)} = 1$, i.e., $\Delta z + h(\Delta x/3) \equiv 0 \pmod 3$. Direct evaluation over all $2 \binom{8}{2} = 56$ same-shore pairs shows that $W_0$ avoids all such congruences for $h \in \{0, 2\}$ and $W_1$ avoids all of them for $h = 1$.
\end{proof}

\subsection{Normal $2$-kernel lifting and theta cover reductions}

\begin{lemma}[Normal $2$-subgroup holonomy lifting]\label{lem:2kernel}
Let $K = \langle a, b \rangle$ be a finite group with $1, a, b$ distinct, and let $N \trianglelefteq K$ be a normal $2$-subgroup. Put $Q = K/N$ and $\bar{x} = xN$.
\begin{enumerate}[label=(\alph*)]
\item If $|\{1, \bar{a}, \bar{b}\}| < 3$, then $\Haar(K; \{1, a, b\})$ contains a simple cycle of length $2^{t+1}$ for some $t \ge 1$.
\item If $|\{1, \bar{a}, \bar{b}\}| = 3$ and $\Haar(Q; \{1, \bar{a}, \bar{b}\})$ contains a simple cycle $C_{2^s}$ ($s \ge 2$), then $\Haar(K; \{1, a, b\})$ contains a simple cycle $C_{2^{s+t}}$ for some $t \ge 0$.
\end{enumerate}
\end{lemma}

\begin{proof}
In (a), two distinct shifts $s, s' \in \{1,a,b\}$ satisfy $s' s^{-1} \in N \setminus \{1\}$, of order $2^t$ ($t \ge 1$). Alternating colors $s$ and $s'$ traces distinct vertices $(0, (s's^{-1})^j)$ and $(1, (s's^{-1})^j s)$ for $0 \le j < 2^t$, closing into a simple $C_{2^{t+1}}$.

In (b), let $w$ be the color word of length $2^s$ tracing the simple cycle in $\Haar(Q; \{1, \bar{a}, \bar{b}\})$. Lifting $w$ from $(0, 1_K)$ in $\Haar(K; \{1,a,b\})$ ends at $(0, h)$ with $h \in N$, of order $2^t$ ($t \ge 0$). Repeating $w$ precisely $2^t$ times gives a cyclically nonbacktracking closed walk of length $2^{s+t}$. Positions inequivalent modulo $2^s$ have distinct images in $\Haar(Q; \{1,\bar{a},\bar{b}\})$, while congruent positions differ by left-multiplication by $h^r$ ($0 < r < 2^t$), which are distinct. Hence the lifted walk is simple.
\end{proof}

\begin{lemma}[Dyadic equal-path theta cover reduction]\label{lem:theta_reduction}
Let $B$ be a finite connected simple graph containing three internally vertex-disjoint paths $\pi_0, \pi_1, \pi_2$ of common length $L = 2^k$ ($k \ge 1$) between two vertices $u, v \in V(B)$. Let $\widetilde{B}$ be a regular $\Gamma$-cover of $B$ for a finite group $\Gamma$. If every component of every simple cubic three-shift Haar graph over every subgroup of $\Gamma$ contains a power-of-two cycle, then every connected component of $\widetilde{B}$ contains a power-of-two cycle.
\end{lemma}

\begin{proof}
Let $g_0, g_1, g_2 \in \Gamma$ be the path voltages of $\pi_0, \pi_1, \pi_2$ from $u$ to $v$. If $g_c = g_{c'}$ for some $c \ne c'$, lifting $\pi_c$ and $\pi_{c'}$ from any vertex in the $u$-fiber produces a simple cycle of length $2L = 2^{k+1}$ in that component of $\widetilde{B}$ by unique path lifting. Otherwise $g_0, g_1, g_2$ are distinct. Relabeling the $v$-fiber by $g_0^{-1}$ realizes the lifted paths as edges of $\Haar(\Gamma; \{1, g_1 g_0^{-1}, g_2 g_0^{-1}\})$. Since $B$ is connected, every component of $\widetilde{B}$ meets the $u$-fiber and contains a component of this Haar core; a simple cycle of length $2^s$ in the Haar core expands along the internally disjoint paths of length $L$ to a simple cycle of length $2^{s+k}$ in $\widetilde{B}$.
\end{proof}

\section{Finite Nilpotent Groups of Class at Most Three}\label{sec:class3}

Throughout Sections~\ref{sec:class3} and~\ref{sec:class4}, we use $[x,y] = x^{-1} y^{-1} x y$ and write the rank-two basic commutators on $a, b$ as
\[
c = [a,b],\qquad d = [c,a],\qquad e = [c,b].
\]
In the free nilpotent group $F_2 / \gamma_4(F_2)$ of rank $2$ and class $3$, every element has a unique ordered Hall normal form $a^x b^y c^z d^u e^v$ with $(x,y,z,u,v) \in \Z^5$. Under the Magnus embedding $a \mapsto 1+X$, $b \mapsto 1+Y$ in $\Z\langle X,Y\rangle/(X,Y)^4$, the leading homogeneous components of $c, d, e$ have unit-determinant leading minors separating all integral normal forms.

\subsection{Universal $C_{32}$ for $p \ge 5$ and for $p = 3$ with $\exp(K_{\mathrm{ab}}) \ge 9$}

Consider the cyclically nonbacktracking $32$-color word
\[
w_{32} = \mathtt{01010102012012101012012010212012}.
\]
Exact computation in $\Z\langle X,Y\rangle/(X,Y)^4$ gives $w_{32} \equiv 1$, and Hall collection yields $(x,y,z,u,v) = (0,0,0,0,0)$. Hence $w_{32}$ closes in \emph{every} two-generator group of nilpotency class at most $3$.

\begin{proposition}[$p \ge 5$ universal $C_{32}$]\label{prop:class3_pge5}
Let $K = \langle a, b \rangle$ have nilpotency class at most $3$ with $[a,b] \ne 1$. Suppose there exists a prime $p \ge 5$ and a homomorphism $\phi : K \to \F_p^2$ sending $a, b$ to a basis of $\F_p^2$. Then $w_{32}$ traces a simple $C_{32}$ in $\Haar(K; \{1,a,b\})$. In particular, every simple cubic Haar graph component over a finite $p$-group of class at most $3$ ($p \ge 5$) contains $C_4, C_8, C_{16}$, or $C_{32}$.
\end{proposition}

\begin{proof}
Across all $2\binom{16}{2} = 240$ same-shore prefix pairs $(i,j)$ ($0 \le i < j < 32$), exact Hall collection of $P_i^{-1} P_j = a^x b^y c^z d^u e^v$ gives abelian gcd $\gcd(|x|, |y|)$ histogram
\[
0:\; 7,\qquad 1:\; 174,\qquad 2:\; 39,\qquad 3:\; 16,\qquad 4:\; 4.
\]
For $p \ge 5$, none of the $233$ pairs with $\gcd(|x|,|y|) \in \{1,2,3,4\}$ can vanish in $\F_p^2$ under $\phi$, so $P_i \ne P_j$ in $K$. The remaining $7$ pairs with $(x,y)=(0,0)$ correspond to $6$-step cyclic arcs:
\begin{align*}
(1,27):&\;\mathtt{120120}, & (2,28):&\;\mathtt{201201}, & (7,13):&\;\mathtt{201201}, & (8,14):&\;\mathtt{012012}, \\
(17,23):&\;\mathtt{012012}, & (18,24):&\;\mathtt{120120}, & (19,25):&\;\mathtt{201201}.
\end{align*}
Since $w_{32}$ closes in $K$, equality of the endpoints of the two wrap-around arcs $(1,27)$ and $(2,28)$ is equivalent to closure of their $6$-step complementary arcs. Each $6$-step arc evaluates to a conjugate of $a^{-1} b a b^{-1}$, $b a b^{-1} a^{-1}$, $b^{-1} a^{-1} b a$, or $a b^{-1} a^{-1} b$, all equal to $1$ if and only if $[a,b] = 1$, contradicting $[a,b] \ne 1$. Thus all $32$ vertices are distinct. For a finite nonabelian $p$-group $K$, Burnside's basis theorem yields $K/\Phi(K) \cong \F_p^2$; if $K$ is abelian, Theorem~\ref{thm:abelian} gives $C_4, C_8$, or $C_{16}$.
\end{proof}

\begin{proposition}[$p = 3$ with $\exp(K_{\mathrm{ab}}) \ge 9$]\label{prop:class3_p3_exp9}
Let $K = \langle a, b \rangle$ be a finite nonabelian $3$-group of nilpotency class at most $3$. If the abelianization $K_{\mathrm{ab}} = K/[K,K]$ has exponent at least $9$, then $\Haar(K; \{1,a,b\})$ contains a simple $C_{32}$.
\end{proposition}

\begin{proof}
Since $K$ is a nonabelian two-generator finite $3$-group, $K_{\mathrm{ab}}/3K_{\mathrm{ab}} \cong K/\Phi(K) \cong \F_3^2$, generated by the images $\alpha, \beta$ of $a, b$ in $K_{\mathrm{ab}}$. If two of $\alpha$, $\beta$, and $\beta - \alpha$ satisfied $3x = 0$, then $3\alpha = 3\beta = 0$, forcing $\exp(K_{\mathrm{ab}}) \le 3$, contrary to hypothesis. Hence at most one of $\alpha, \beta, \beta - \alpha$ is killed by $3$. Re-anchoring the shift set $\{1,a,b\}$ at an appropriate $t \in \{1,a,b\}$ replaces $(a,b)$ by $(a,b)$, $(a^{-1}, a^{-1}b)$, or $(b^{-1}, b^{-1}a)$, giving an isomorphic Haar graph generating $K$ whose new generators both satisfy $3\bar{a} \ne 0$ and $3\bar{b} \ne 0$ in $K_{\mathrm{ab}}$.

Evaluate $w_{32}$ on this re-anchored pair. The $217$ same-shore prefix pairs with $\gcd(|x|,|y|) \in \{1,2,4\}$ are nonzero modulo $3K_{\mathrm{ab}}$. The $16$ pairs with $\gcd(|x|,|y|) = 3$ satisfy
\[
(x,y) \in \{(-3,0)^5,\; (3,0)^9,\; (0,3)^2\},
\]
giving $\pm 3\bar{a}$ or $3\bar{b}$, nonzero in $K_{\mathrm{ab}}$. The $7$ pairs with $(x,y)=(0,0)$ force $[a,b]=1$ as in Proposition~\ref{prop:class3_pge5}, impossible since $K$ is nonabelian. Thus $w_{32}$ is a simple $C_{32}$.
\end{proof}

\subsection{The remaining $p = 3$, $K_{\mathrm{ab}} \cong C_3 \times C_3$ branch}

\begin{lemma}[Cube centrality in class-$3$ $3$-groups]\label{lem:p3_cubes}
Let $K = \langle a, b \rangle$ be a finite nonabelian $3$-group of nilpotency class exactly $3$ with $K_{\mathrm{ab}} \cong C_3 \times C_3$. Then $a^3, b^3 \in \gamma_3(K) \le Z(K)$ and $c^3 = 1$, while $c \notin \gamma_3(K)$.
\end{lemma}

\begin{proof}
In class $3$, $\gamma_3(K) = \langle d, e \rangle \le Z(K)$ and $\gamma_2(K) = \langle c, d, e \rangle$ is abelian. Because $K_{\mathrm{ab}} \cong C_3^2$, we have $a^3, b^3 \in \gamma_2(K)$. In $F_2/\gamma_4(F_2)$, exact commutator collection gives
\[
[a^3, b] = c^3 d^3, \qquad [a, b^3] = c^3 e^3.
\]
Because $a^3, b^3 \in \gamma_2(K)$, both left-hand sides lie in $\gamma_3(K)$, so $c^3 \in \gamma_3(K) \le Z(K)$. Centrality gives $d^3 = [c,a]^3 = [c^3, a] = 1$ and $e^3 = [c^3, b] = 1$.

Write $a^3 = c^r d^s e^t$ and $b^3 = c^u d^v e^w$ with $r, u \in \{0,1,2\}$ (absorbing any $c^3$ into $\gamma_3(K)$). From $[a^3, a] = 1$ and $[b^3, b] = 1$ we obtain $d^r = 1$ and $e^u = 1$. On the other hand, substituting into the commutators gives $[a^3, b] = [c^r, b] = e^r$ and $[a, b^3] = [a, c^u] = d^{-u}$, whence $c^3 = e^r = d^{-u}$. If $r \not\equiv 0 \pmod 3$, then $d^r = 1$ and $d^3 = 1$ imply $d = 1$; hence $c^3 = d^{-u} = 1$, forcing $e^r = 1$ and $e = 1$, contradicting $\gamma_3(K) \ne 1$. Thus $r = 0$, and symmetrically $u = 0$. Therefore $a^3, b^3 \in \gamma_3(K) \le Z(K)$ and $c^3 = 1$. Finally, if $c \in \gamma_3(K)$, then $\gamma_2(K) = \gamma_3(K)$ would be central, forcing $\gamma_3(K) = 1$, a contradiction.
\end{proof}

\begin{proposition}[Balanced $C_{32}$ for $K_{\mathrm{ab}} \cong C_3^2$]\label{prop:class3_p3_c3xc3}
Let $K = \langle a, b \rangle$ be a finite nonabelian $3$-group of class exactly $3$ with $K_{\mathrm{ab}} \cong C_3 \times C_3$. Then $\Haar(K; \{1,a,b\})$ contains a simple $C_{32}$ whose color word is balanced (zero total signed exponent in both $a$ and $b$). Specifically, if $e = [[a,b],b] \ne 1$, the word
\[
w_{32}^{\mathrm{bal}} = \mathtt{01010201021201201020201210212012}
\]
traces a simple $C_{32}$; if $e = 1$, swapping colors $1$ and $2$ traces a simple $C_{32}$.
\end{proposition}

\begin{proof}
Exact Magnus expansion in $\Z\langle X,Y\rangle/(X,Y)^4$ verifies that $w_{32}^{\mathrm{bal}}$ is cyclically nonbacktracking, evaluates to $1$ in $F_2/\gamma_4(F_2)$, and has total signed exponents $(0,0)$. Across all $240$ same-shore prefix pairs $P_i^{-1} P_j = a^x b^y c^z d^u e^v$, exact Hall collection reveals that \emph{only five pairs} satisfy $3 \mid x$, $3 \mid y$, and $3 \mid z$, namely
\[
(i,j) \in \{(3,23),\; (4,22),\; (5,21),\; (6,20),\; (7,19)\},
\]
and all five have exact integer Hall coordinates $(x,y,z,u,v) = (0,0,0,0,-1)$, i.e., $P_i^{-1} P_j = e^{-1}$.

By Lemma~\ref{lem:p3_cubes}, $a^3, b^3 \in \gamma_3(K)$ and $c \gamma_3(K)$ has order $3$ in $\gamma_2(K)/\gamma_3(K)$. Hence $P_i \equiv P_j \pmod{\gamma_3(K)}$ if and only if $x \equiv y \equiv z \equiv 0 \pmod 3$. Thus only these five pairs can collide in $K$, and each yields $e^{-1}$. If $e \ne 1$, no pair collides and $w_{32}^{\mathrm{bal}}$ is a simple balanced $C_{32}$. If $e = 1$, class $3$ forces $d \ne 1$; swapping $a$ and $b$ gives $c' = [b,a] = c^{-1}$ and $e' = [c', a] = [c^{-1}, a] = d^{-1} \ne 1$ (by centrality of $d$), so the color-swapped balanced word yields a simple $C_{32}$.
\end{proof}

\subsection{Proof of Theorem~\ref{thm:class3}}

\begin{proof}[Proof of Theorem~\ref{thm:class3}]
Normalize the component shifts to $\{1, a, b\}$ inside $K = \langle a, b \rangle \le \Gamma$. Then $N_K = K \cap N$ is a normal $2$-subgroup of $K$ whose quotient $Q = K/N_K$ embeds into $\Gamma/N$, hence is finite nilpotent of class at most $3$. By Lemma~\ref{lem:2kernel}, it suffices to exhibit a simple power-of-two cycle in $\Haar(Q; \{1, \bar{a}, \bar{b}\})$ whenever $1, \bar{a}, \bar{b}$ remain distinct in $Q$.

If $Q$ is abelian, Theorem~\ref{thm:abelian} gives $C_4, C_8$, or $C_{16}$. Otherwise $Q$ is nonabelian of odd order and nilpotency class $2$ or $3$. Decompose $Q = \prod_p Q_p$ into its odd-order Sylow $p$-subgroups.
\begin{itemize}
\item If some Sylow subgroup $Q_p$ with $p \ge 5$ requires two generators, the projections of $\bar{a}, \bar{b}$ form a basis of $Q_p/\Phi(Q_p) \cong \F_p^2$, and Proposition~\ref{prop:class3_pge5} directly yields a simple $C_{32}$ in $Q$.
\item Otherwise every Sylow $p$-subgroup for $p \ge 5$ is cyclic, so $Q = P_3 \times C_m$ with $\gcd(m,6) = 1$, $C_m$ cyclic central, and $P_3$ nonabelian of class $2$ or $3$ (generated by the $P_3$-projections $\tilde{a}, \tilde{b}$ of $\bar{a}, \bar{b}$).
\begin{itemize}
\item If $P_3$ has class $2$, then $Q$ itself has class $2$ and odd order, so $\ord([\bar{a}, \bar{b}]) \ge 3$; Proposition~\ref{prop:class2} gives a simple $C_{16}$ in $Q$.
\item If $P_3$ has class $3$ and $\exp((P_3)_{\mathrm{ab}}) \ge 9$, then $\exp(Q_{\mathrm{ab}}) \ge 9$ and Proposition~\ref{prop:class3_p3_exp9} (applied via $Q \twoheadrightarrow P_3$, or directly noting that the re-anchored generators of $Q$ inherit nontrivial cubes in $P_3/[P_3,P_3]$ and nonabelian class $3$ in $Q$) gives a simple $C_{32}$ in $Q$.
\item If $P_3$ has class $3$ and $(P_3)_{\mathrm{ab}} \cong C_3 \times C_3$, Proposition~\ref{prop:class3_p3_c3xc3} provides a \emph{balanced} simple $C_{32}$ color word in $\Haar(P_3; \{1, \tilde{a}, \tilde{b}\})$. Because its signed integer exponents in both generators are zero, this same word closes in the abelian factor $C_m$, and its simple projection in $P_3$ forces all $32$ vertices in $Q = P_3 \times C_m$ to be distinct, giving a simple $C_{32}$ in $Q$.
\end{itemize}
\end{itemize}
Notice that when $\Gamma$ itself is odd-order nilpotent of class at most $3$, $N=1$, so $Q=\Gamma$ already contains a simple $C_4, C_8, C_{16}$, or $C_{32}$.
\end{proof}

\section{Unconditional Class-Four Theorem, Rigidity, and Boundary Families}\label{sec:class4}

In the free rank-two nilpotent group $F_2/\gamma_5(F_2)$ of class $4$, write the basic commutators on $a, b$ as $c = [a,b]$, $d = [c,a]$, $e = [c,b]$, and
\[
f = [d,a],\qquad g = [d,b],\qquad h = [e,b].
\]
By the Hall--Witt identity modulo $\gamma_5(F_2)$, $[e,a] = [[c,b],a] = [[c,a],b] = g$ since $\gamma_4$ is central and $[b,a] = c^{-1}$ commutes with $c$ modulo $\gamma_3$. Moreover, $\gamma_2(G)$ is abelian in every class-$4$ group ($\gamma_2$ commutes with $\gamma_2$, leaving $[\gamma_2, \gamma_3] \le \gamma_5 = 1$). Every element of $F_2/\gamma_5(F_2)$ has unique integral Hall normal form
\[
a^x b^y c^z d^u e^v f^r g^s h^t \qquad ((x,y,z,u,v,r,s,t) \in \Z^8).
\]
Under $a \mapsto 1+X, b \mapsto 1+Y$ in $\Z\langle X,Y\rangle/(X,Y)^5$, the degree-$4$ leading matrix on $(X^3 Y, X^2 Y^2, X Y^3)$ versus $(f, g, h)$ has determinant $-1$, certifying faithful integral collection. Exact computation in $F_2/\gamma_5(F_2)$ also yields the cube commutator identities
\begin{equation}\label{eq:cube_comm_4}
[a^3, b] = c^3 d^3 f,\quad [a, b^3] = c^3 e^3 h,\quad [c^3, a] = d^3,\quad [c^3, b] = e^3.
\end{equation}

\subsection{Sharp $C_{32}$ for $p \ge 5$ and two-stage descent for $p = 3$}

Define two balanced cyclically nonbacktracking $32$-color words:
\begin{align*}
w_4 &= \mathtt{01012012102021010202102120101202}, \\
w_5 &= \mathtt{01012020101201210202101020210212}.
\end{align*}
Both words evaluate to the identity $(0)^8$ in $F_2/\gamma_5(F_2)$. Across their $240$ same-shore prefix differences $P_i^{-1} P_j = a^x b^y c^z d^u e^v f^r g^s h^t$, both $w_4$ and $w_5$ have abelian gcd $\gcd(|x|, |y|)$ histogram $\{0:19,\ 1:177,\ 2:44\}$ (hence zero pairs with nonzero $(x,y) \equiv (0,0) \pmod 3$), and all $19$ pairs with $(x,y)=(0,0)$ satisfy $z \in \{-1, 1, 3, 4\}$.

\begin{proposition}[Sharp $C_{32}$ for all class-$\le 4$ $p$-groups with $p \ge 5$]\label{prop:class4_pge5_c32}
Let $p \ge 5$ be a prime and let $G = \langle a, b \rangle$ be a nonabelian finite $p$-group of nilpotency class at most $4$ (arbitrary exponent). Then $w_4$ traces a simple balanced $C_{32}$ in $\Haar(G; \{1,a,b\})$.
\end{proposition}

\begin{proof}
Since $G$ is nonabelian of class $\le 4$ and $2$-generated, $G/\Phi(G) \cong \F_p^2$ on the images of $a, b$, and $c = [a,b] \notin \gamma_3(G)$ has order divisible by $p$ modulo $\gamma_3(G)$. Because $p \ge 5$, none of the $221$ same-shore differences of $w_4$ with $\gcd(|x|,|y|) \in \{1, 2\}$ can vanish in $G/\Phi(G)$, and none of the $19$ differences with $(x,y)=(0,0)$ and $z \in \{-1,1,3,4\}$ can vanish in $\gamma_2(G)/\gamma_3(G)$.
\end{proof}

We now reduce finite class-$4$ $3$-groups $G = \langle a, b \rangle$. If $G$ has class $\le 3$, Propositions~\ref{prop:class2}, \ref{prop:class3_p3_exp9}, and~\ref{prop:class3_p3_c3xc3} give $C_{16}$ or balanced $C_{32}$. Assume $G$ has class exactly $4$, so $\gamma_4(G) \ne 1$, $c \notin \gamma_3(G)$, and $G/\gamma_4(G)$ has class $3$.

\begin{lemma}[Two-stage elimination at $p = 3$]\label{lem:p3_descent}
Let $G = \langle a, b \rangle$ be a finite $3$-group of nilpotency class $4$.
\begin{enumerate}[label=(\arabic*)]
\item If $\ord(c \bmod \gamma_3(G)) \ge 9$, then $w_4$ traces a simple balanced $C_{32}$ in $\Haar(G; \{1,a,b\})$.
\item Otherwise $c^3 \in \gamma_3(G)$. Then at least one of $w_4, w_5$ traces a simple balanced $C_{32}$ in $\Haar(G; \{1,a,b\})$ unless BOTH
\[
c^3 \in \gamma_3(G)^3 \gamma_4(G) \quad\text{and}\quad d e \in \gamma_3(G)^3 \gamma_4(G).
\]
\item When $c^3, de \in \gamma_3(G)^3 \gamma_4(G)$, we have $\gamma_3(G)^3 \le \gamma_4(G)^3$ and $\gamma_4(G)/\gamma_4(G)^3 \cong \F_3$ generated by $\bar{f}$. Modulo $\gamma_4(G)^3$, writing $c^3 = \bar{f}^\mu$ and $de = \bar{f}^\lambda$ defines the quotient group $Q_{\mathrm{crit}}(\mu, \lambda)$ ($\mu, \lambda \in \F_3$), and $\{w_4, w_5\}$ leave uncovered only $(\mu, \lambda) = (2, 1)$.
\end{enumerate}
\end{lemma}

\begin{proof}
For~(1), if $\ord(c \bmod \gamma_3(G)) \ge 9$, then $c^z \notin \gamma_3(G)$ for all $z \in \{-1,1,3,4\}$, and all remaining same-shore differences of $w_4$ have $\gcd(|x|,|y|) \in \{1,2\}$ (nonzero in $G/\Phi(G) \cong \F_3^2$).

For~(2), suppose $c^3 \in \gamma_3(G)$, and set $V = \gamma_3(G)/(\gamma_3(G)^3 \gamma_4(G)) \ne 0$. Notice $\gamma_3(G)/\gamma_4(G)$ is abelian and generated by $\bar{d}, \bar{e}$. For $w_4$, every same-shore difference with $(x,y)=(0,0)$ and $3 \mid z$ satisfies $(z, u \bmod 3, v \bmod 3) = (3, 1, 1)$; for $w_5$, every such difference satisfies $(3, 2, 2)$. For any linear functional $\varphi : V \to \F_3$, let $(D, E, C_3) = (\varphi(\bar{d}), \varphi(\bar{e}), \varphi(c^3))$. Testing all $24$ triples in $(\F_3^2 \setminus \{(0,0)\}) \times \F_3$ eliminates $22/24$, leaving only $(D, E, C_3) \in \{(1,2,0), (2,1,0)\}$. Thus unless $c^3$ and $de$ vanish identically in $V$, some $\varphi$ detects a nonzero collision obstruction in $w_4$ or $w_5$.

For~(3), writing $c^3 = d^{3p} e^{3q} w_c$ ($w_c \in \gamma_4(G)$) and commuting with $a, b$ gives $d^3 = [c^3, a] = (f^p g^q)^3 \in \gamma_4(G)^3$ and $e^3 = [c^3, b] = (g^p h^q)^3 \in \gamma_4(G)^3$, so $\gamma_3(G)^3 \le \gamma_4(G)^3$. Hence $de \in \gamma_4(G) \bmod \gamma_4(G)^3$, and commuting with $a, b$ gives $fg, gh \in \gamma_4(G)^3$. Since $\gamma_4(G) \ne 1$ is abelian, Burnside's basis theorem yields $W = \gamma_4(G)/\gamma_4(G)^3 \cong \F_3 \cdot \bar{f}$ with $\bar{g} = -\bar{f}$ and $\bar{h} = \bar{f}$. Passing to $G/\gamma_4(G)^3$, $c^3 = \bar{f}^\mu$ and $de = \bar{f}^\lambda$ determine the $5$-coordinate quotient $Q_{\mathrm{crit}}(\mu,\lambda)$ of $F_2/\gamma_5(F_2)$ via $z = 3q_z + z_0$, $u \equiv u_0 + v \pmod 3$, and $w \equiv \mu q_z + \lambda v + r - s + t \pmod 3$. Evaluating all $9$ pairs $(\mu, \lambda) \in \F_3^2$ eliminates $8/9$ pairs via $w_4$ or $w_5$, leaving uniquely $(\mu, \lambda) = (2, 1)$.
\end{proof}

Consider next the two balanced cyclically nonbacktracking $32$-color words
\[
w_0 = \mathtt{01010120121010210201210210102012},\qquad
w_0^{\mathrm{swap}} = \mathtt{02020210212020120102120120201021}.
\]
Both close identically in $F_2/\gamma_5(F_2)$. Their nonzero $(x,y) \equiv (0,0) \pmod 3$ differences occur solely at $(x,y) = (\pm 3, 0)$ (for $w_0$) and $(0, \pm 3)$ (for $w_0^{\mathrm{swap}}$), and both induce zero collisions among $(x,y)=(0,0)$ pairs inside $Q_{\mathrm{crit}}(2, 1)$. Consequently, if $\exp(G_{\mathrm{ab}}) \ge 9$, then at least one of $3\bar{a} \ne 0$ (where $w_0$ is simple) or $3\bar{b} \ne 0$ (where $w_0^{\mathrm{swap}}$ is simple) holds in $G_{\mathrm{ab}}$, resolving every $\exp(G_{\mathrm{ab}}) \ge 9$ branch unconditionally.

\subsection{Algebraic Rigidity Lemma for $G_{\mathrm{ab}} \cong C_3^2$ and proofs of Theorems~\ref{thm:class4} and~\ref{thm:theta}}

\begin{lemma}[Algebraic rigidity at $G_{\mathrm{ab}} \cong C_3 \times C_3$]\label{lem:rigidity_crit}
Suppose $G = \langle a, b \rangle$ is a finite $3$-group of nilpotency class $4$ with $G_{\mathrm{ab}} \cong C_3 \times C_3$ lying in the remaining branch of Lemma~\ref{lem:p3_descent} ($c^3, de \in \gamma_3(G)^3 \gamma_4(G)$ and $(\mu,\lambda)=(2,1)$). Then $f^3 = 1$ identically in $G$, $|G| = 243$, and $G \cong G_{\mathrm{crit}}(\alpha, \beta)$ ($a^3 = f^\alpha$, $b^3 = f^\beta$) for some $(\alpha, \beta) \in \F_3^2$. Moreover:
\begin{enumerate}[label=(\alph*)]
\item For all $8$ pairs $(\alpha,\beta) \ne (0,0)$, $w_0$ ($\alpha \ne 0$) or $w_0^{\mathrm{swap}}$ ($\alpha=0, \beta \ne 0$) yields a simple balanced $C_{32}$.
\item For the unique group $G_{\mathrm{crit}}(0,0)$, none of the $288$ free-class-$4$ $32$-words survive, yet the balanced non-free word $w_{32,\mathrm{crit}}^{\mathrm{bal}} = \mathtt{02020101210120101021021012101021}$ yields a simple balanced $C_{32}$, and the free-class-$4$ balanced $64$-word
\[
w_{64,\mathrm{crit}} = \mathtt{0101210120120212102012020120102021012012120121202102121012021021}
\]
traces a simple balanced $C_{64}$ across all $9$ groups $G_{\mathrm{crit}}(\alpha,\beta)$.
\end{enumerate}
\end{lemma}

\begin{proof}
Applying Lemma~\ref{lem:p3_cubes} to the class-$3$ quotient $G/\gamma_4(G)$ shows $c^3 \in \gamma_4(G) \le Z(G)$ and $a^3, b^3 \in \gamma_3(G)$ in $G$ itself. Commuting $c^3 \in Z(G)$ with $a,b$ gives $d^3 = e^3 = 1$ identically in $G$, whence $\gamma_3(G)^3 = \gamma_4(G)^3$ (since $\gamma_3(G)$ is abelian). Next, $de \in \gamma_3(G)^3 \gamma_4(G) = \gamma_4(G) \le Z(G)$ gives $fg = [de,a] = 1$ and $gh = [de,b] = 1$ in $G$, so $\gamma_4(G) = \langle f \rangle$ is cyclic. Write $a^3 = d^s e^t w_a$ and $b^3 = d^u e^v w_b$ with $s,t,u,v \in \{0,1,2\}$ and $w_a, w_b \in \gamma_4(G)$. Then $1 = [a^3,a] = f^{s-t}$ and $1 = [b^3,b] = f^{v-u}$ in $G$ force $s=t$ and $u=v$ (because $\ord(f) \ge 3$ is a power of $3$ and $|s-t|,|v-u| \le 2$). Since $de \in \gamma_4(G)$, this puts $a^3, b^3 \in \gamma_4(G) = \langle f \rangle$. Finally, writing $de = f^{\lambda_0}$ with $\lambda_0 \equiv 1 \pmod 3$, centrality of $d^3=1, e^3=1$ gives $f^{3\lambda_0} = (de)^3 = d^3 e^3 = 1$, forcing $f^3 = 1$ identically in $G$ (no order-$729$ extension exists). Hence $|G| = 3^5 = 243$ and $G \cong G_{\mathrm{crit}}(\alpha,\beta)$ with $a^3=f^\alpha, b^3=f^\beta$. Parts~(a) and~(b) are verified exhaustively across all $9$ isomorphism types and all $288$ free-class-$4$ $32$-words.
\end{proof}

\begin{proof}[Proof of Theorems~\ref{thm:class4} and~\ref{thm:theta} (for Class $\le 4$)]
Normalize a component of $\Haar(\Gamma; S)$ to $K = \langle a, b \rangle \le \Gamma$ and pass to $Q = K/(K \cap N) = \prod_p Q_p$ via Lemma~\ref{lem:2kernel}. If two shifts coincide in $Q$, Lemma~\ref{lem:2kernel}(a) gives a power-of-two cycle upstairs. Otherwise $Q$ is finite nilpotent of class at most $4$. If $Q$ is abelian or has class $\le 3$, Theorems~\ref{thm:abelian} and~\ref{thm:class3} yield $C_4, C_8, C_{16}$, or $C_{32}$ in $Q$, which lifts via Lemma~\ref{lem:2kernel}(b).

Suppose $Q$ has class $4$. If some Sylow $p$-subgroup $Q_p$ ($p \ge 5$) has class $4$ (or class $2,3$ nonabelian), Proposition~\ref{prop:class4_pge5_c32} (or Proposition~\ref{prop:class2}/Proposition~\ref{prop:class3_pge5}) provides a balanced simple $C_{16}$ or $C_{32}$ word closing identically in $F_2/\gamma_5(F_2)$; since $\cl(Q) \le 4$, this word closes in $Q$ and separates all vertices via $Q \twoheadrightarrow Q_p$. Otherwise every $Q_p$ ($p \ge 5$) is cyclic, so $Q = P_3 \times C_m$ with $\gcd(m,6)=1$, $C_m$ cyclic central, and $P_3$ of class $4$. By Lemmas~\ref{lem:p3_descent} and~\ref{lem:rigidity_crit}, $\Haar(P_3; \{1,\tilde{a},\tilde{b}\})$ always contains a \emph{balanced} simple $C_{32}$ word ($w_4, w_5, w_0, w_0^{\mathrm{swap}}$, or $w_{32,\mathrm{crit}}^{\mathrm{bal}}$). Because $w_{32,\mathrm{crit}}^{\mathrm{bal}}$ closes in $G_{\mathrm{crit}}(0,0)$ and has $\Z^2$-balance $(0,0)$, each such balanced $32$-word closes identically in $P_3 \times C_m$ and keeps all $32$ vertices distinct via $Q \twoheadrightarrow P_3$, yielding a simple $C_{32}$ in $Q$, which lifts via Lemma~\ref{lem:2kernel}(b) to $\Haar(K; \{1,a,b\})$. Notice that when $\Gamma$ itself is odd-order nilpotent of class $\le 4$, $N = 1$ and $Q = \Gamma$, establishing the sharp $C_{32}$ bound. Theorem~\ref{thm:theta} follows directly from Lemma~\ref{lem:theta_reduction} once Theorem~\ref{thm:class6} is established below.
\end{proof}

\subsection{Degree-$6$ Magnus calculus and proof of Theorem~\ref{thm:class6} and Propositions~\ref{prop:boundaries}--\ref{prop:barrier32}}

For $p \ge 5$, we exploit the exact degree-$6$ Magnus closure of three companion $64$-color words obtaining central determinant $27$:
\begin{align*}
w_{64}^{(1)} &= \mathtt{0101012012021020101021021012012101020120210210101012102101201201}, \\
w_{64}^{(2)} &= \mathtt{0202021021012010202012012021021202010210120120202021201202102102}, \\
w_{64}^{(3)} &= \mathtt{1010102102120121010120120102102010121021201201010102012010210210}.
\end{align*}

\begin{lemma}[Degree-$6$ Magnus vanishing, all-prime gcd $1$, and prefix structure]\label{lem:c64_prefixes}
Each word $w_{64}^{(k)}$ ($k=1,2,3$) evaluates identically to $1$ in $\Z\langle X,Y\rangle/(X,Y)^7$ (leaving $67, 67, 20$ nonzero terms at degree $7$), hence belongs to $D_7(F_2) = \gamma_7(F_2)$ and closes identically in every two-generator group of nilpotency class at most $6$. Moreover:
\begin{enumerate}[label=(\alph*)]
\item Across all $992$ same-shore prefix differences of $w_{64}^{(1)}$ in $\Z\langle X,Y\rangle/(X,Y)^7$, \emph{every single difference} (`992/992`) has integral monomial coefficient gcd equal to $1$.
\item Projecting to $F_2/\gamma_5(F_2)$, each $w_{64}^{(k)}$ satisfies: (1) $80$ pairs have $(x,y) = (0,0)$ and $912$ satisfy $1 \le \gcd(|x|, |y|) \le 4$ ($1:668,\ 2:136,\ 3:84,\ 4:24$); (2) $28$ zero-$(x,y)$ pairs have $1 \le |z| \le 4$; and (3) the remaining $52$ pairs split into $47$ grade-$3$ pairs $(u,v) = \lambda_k (u_0^{(k)}, v_0^{(k)})$ ($|\lambda_k| \in \{1,2\}$, directions $(1,-1), (1,-1), (0,1)$) and $5$ central pairs $\pm v_k$ where $v_1 = (-3, 3, 0)$, $v_2 = (0, -3, 3)$, $v_3 = (0, 3, 0)$ ($\det = 27$).
\end{enumerate}
\end{lemma}

\begin{proof}[Proof of Theorem~\ref{thm:class6}]
Normalize to $K = \langle a, b \rangle \le \Gamma$ and pass modulo $K \cap N$ to $Q = \prod_{p \ge 3} Q_p$ of nilpotency class at most $6$ with $\cl(Q_3) \le 4$. If shifts collide in $Q$, Lemma~\ref{lem:2kernel}(a) finishes immediately; if all $Q_p$ ($p \ge 5$) have class $\le 4$, then $\cl(Q) \le 4$ and Theorem~\ref{thm:class4} produces $C_4, C_8, C_{16}$, or $C_{32}$ in $Q$.

Assume instead that some prime $p^* \ge 5$ has $\cl(Q_{p^*}) \in \{5, 6\}$, and set $\overline{Q}_{p^*} = Q_{p^*}/\gamma_5(Q_{p^*})$. If $\gamma_4(\overline{Q}_{p^*}) = 1$, then $\gamma_4(Q_{p^*}) \le \gamma_5(Q_{p^*})$ would force $\gamma_4(Q_{p^*}) = 1$ by nilpotency, contradicting $\cl(Q_{p^*}) \ge 5$. Hence $\overline{Q}_{p^*}$ has class exactly $4$, cannot be cyclic, and since $\gamma_5(Q_{p^*}) \le [Q_{p^*}, Q_{p^*}] \le \Phi(Q_{p^*})$, Burnside's basis theorem gives $\overline{Q}_{p^*}/\Phi(\overline{Q}_{p^*}) \cong \F_{p^*}^2$ on $\{\bar{a}_{p^*}, \bar{b}_{p^*}\}$, keeping the three shifts distinct in $\overline{Q}_{p^*}$.

Because $\cl(Q) \le 6$, Lemma~\ref{lem:c64_prefixes} ensures that all three words $w_{64}^{(1)}, w_{64}^{(2)}, w_{64}^{(3)}$ close identically in $Q$. It suffices to exhibit one $w_{64}^{(k)}$ having zero same-shore collisions inside $\overline{Q}_{p^*}$:
\begin{enumerate}[label=(\roman*)]
\item Since $p^* \ge 5$, none of the $912$ pairs with $1 \le \gcd(|x|,|y|) \le 4$ vanish in $\overline{Q}_{p^*}/\Phi(\overline{Q}_{p^*}) \cong \F_{p^*}^2$, and none of the $28$ pairs with $x=y=0, 1 \le |z| \le 4$ vanish modulo $\gamma_3(\overline{Q}_{p^*})$ (as $\gamma_2/\gamma_3 = \langle \bar{c}\rangle$ is cyclic of $p^*$-power order $\ge 5$).
\item Across the $47$ grade-$3$ pairs, since $\gcd(|\lambda_k|, p^*) = 1$, a collision occurs iff $d e^{-1} \in \gamma_4(\overline{Q}_{p^*})$ (for $k=1,2$) or $e \in \gamma_4(\overline{Q}_{p^*})$ (for $k=3$).
\item If $d e^{-1} \in \gamma_4(\overline{Q}_{p^*})$, writing $d = e w$ ($w \in Z(\overline{Q}_{p^*})$) forces $f = [d,a] = [e,a] = g$ and $g = [d,b] = [e,b] = h$. Since $\gamma_4(\overline{Q}_{p^*}) = \langle f,g,h\rangle \ne 1$, we have $g \ne 1$ (hence $g^{\pm 3} \ne 1$ for $p^* \ge 5$) and $e \notin \gamma_4(\overline{Q}_{p^*})$ (else $h = 1 \implies \gamma_4 = 1$). Thus $w_{64}^{(3)}$ is collision-free in $\overline{Q}_{p^*}$.
\item If $e \in \gamma_4(\overline{Q}_{p^*}) \le Z(\overline{Q}_{p^*})$, then $g = h = 1$, forcing $f \ne 1$ and $d \notin \gamma_4(\overline{Q}_{p^*})$. Hence $(d e^{-1})^{\lambda_1} \equiv d^{\lambda_1} \not\equiv 1 \pmod{\gamma_4}$ and $f^{-3} g^3 = f^{-3} \ne 1$, so $w_{64}^{(1)}$ is collision-free in $\overline{Q}_{p^*}$.
\item Otherwise neither $d e^{-1}$ nor $e$ lies in $\gamma_4(\overline{Q}_{p^*})$, so zero grade-$3$ collisions occur in any $w_{64}^{(k)}$. Since $p^* \ge 5$ and $\gamma_4(\overline{Q}_{p^*}) \ne 1$, the system $f^{-3} g^3 = g^{-3} h^3 = g^3 = 1$ admits only $f = g = h = 1$, impossible. Thus at least one $w_{64}^{(k)}$ traces a simple $C_{64}$ in $\overline{Q}_{p^*}$, hence in $Q$, and lifts via Lemma~\ref{lem:2kernel}(b) to $\Haar(K; \{1,a,b\})$. Direct matrix checks confirm all three $w_{64}^{(k)}$ trace simple $C_{64}$ cycles in $UT_7(\F_5)$ (order $5^{21}$, class $6$, $A^5 \ne I = A^{25}$) and in $G_{\mathrm{crit}}(0,0) \times UT_7(\F_5)$.
\end{enumerate}
\end{proof}

\begin{proof}[Proof of Propositions~\ref{prop:boundaries} and~\ref{prop:barrier32}]
For Proposition~\ref{prop:boundaries}(a), let $A^{(6)} = \Z\langle X,Y\rangle/(X,Y)^7$ and $K_{p,k}^{(6)} = \langle 1+X, 1+Y \rangle \le (A^{(6)}/p^k A^{(6)})^\times$. By Lemma~\ref{lem:c64_prefixes}(a), $w_{64}^{(1)}$ closes identically to $1$ in $A^{(6)}$ and all $992$ same-shore prefix differences have integral coefficient gcd $1$ (`1:992`). Thus no difference vanishes modulo any prime $p \ge 2$, nor modulo $p^k$ ($k \ge 1$). Literal truncated-algebra walks for $p \in \{2,3,5\}$ modulo $p$ and $p^2$ ($\ord(1+X) = 16, 27, 125$) confirm simple $C_{64}$ cycles throughout.

For part~(b), in $UT_4(\F_p)$ with $A = I + E_{12} + E_{23}$ and $B = I + E_{23} + E_{34}$, exact commutator arithmetic gives $[A,B] = I + E_{13} - E_{14} + E_{24}$, $D = [[A,B], A] = I - E_{14}$ (under $xyx^{-1}y^{-1}$), and $E = [[A,B], B] = D^{-1}$, proving class $3$ and order $p^4$ for all $p \ge 7$. Exhaustive evaluation of all $71{,}004$ cyclically nonbacktracking words of lengths $4, 6, 8, 10, 12, 16$ over $\Z$ yields nonzero entry-gcds whose prime factors are at most $5$, so no cycle of length $\le 16$ exists modulo $p \ge 7$ except odd lengths or $C_{14}$; the explicit integral words $\mathtt{01010210202012}$ (all $42$ prefix-difference gcds equal $1$) and $\mathtt{02101020201010120210210101010101}$ (all $240$ prefix-difference gcds in $\{1,2,3,5\}$) give girth $14$ and a simple $C_{32}$ for every $p \ge 7$.

For part~(c), in $UT_6(\F_p)/\langle I + E_{16}\rangle$ with $A = I + E_{12} + E_{34} + E_{56}$ and $B = I + E_{23} + E_{45}$, exact computation gives $G = [[[A,B],B], A] = I + 2 E_{15} - 2 E_{26}$ of order $p \ge 11$, class $4$, and $|K| = p^6$. Exhaustive evaluation of all $87{,}390$ cyclically nonbacktracking words of lengths $4, 6, \dots, 16$ over $\Z$ yields nonzero entry-gcds with prime factors only in $\{2,3,5,7\}$, proving girth $\ge 18$ modulo $p \ge 11$. Meanwhile $\mathtt{01010101201210210210102102102012}$ (240 prefix-difference gcds in $\{1,2,3,4\}$) and $w_{\mathrm{p3}} = \mathtt{0101010102020120101020121020201202102020121020101021020201010101}$ (992 prefix-difference gcds in $\{1,\dots,8\}$) yield simple $C_{32}$ and $C_{64}$ for all $p \ge 11$.

Finally, Proposition~\ref{prop:barrier32} follows directly from evaluating all $288$ free-class-$4$ $32$-color words in $\Z\langle X,Y\rangle/(X,Y)^6$, where $0/288$ close at degree $5$.
\end{proof}

\section{Open Frontier, Reproducibility, and Provenance}

With nilpotency class at most four unconditionally resolved across all primes and exponents, and classes five and six resolved for all $p \ge 5$ (arbitrary exponent) and class-$\le 4$ Sylow-$3$ quotients, two natural algebraic frontiers remain open for cubic Haar graphs:
\begin{enumerate}
\item $3$-groups of nilpotency class $5$ or $6$ (and arbitrary nilpotency class $c \ge 7$), where Proposition~\ref{prop:barrier32} confirms no universal free-group $32$-word exists for $c \ge 5$ and an inductive structural quotient-lift (adapted to cubic subgraphs such as equal-length dyadic thetas) appears preferable to searching for higher-class universal words; and
\item General non-nilpotent solvable cubic Haar extensions (with all semidirect $C_2$-, $C_3$-, and $C_4$-extensions resolved in companion Papers~III and~IV).
\end{enumerate}

\subsection*{Reproducibility and AI Disclosure}
Every word identity, prefix-difference gcd histogram, Hall coordinate table, control group order, and cycle witness in this paper is independently reproduced in $\sim 1.5$ seconds by the zero-dependency Python~3 script \texttt{paper1\_nilpotent\_haar/verify\_paper1.py}, with full output archived in \texttt{paper1\_nilpotent\_haar/certificates/paper1\_certificates.json}. Portions of the symbolic discovery, algebraic drafting, and verification scripts were developed with AI assistance; all mathematical theorems, proofs, and executable verifier assertions were audited and stand on the accompanying deterministic code and certificates.

\begin{thebibliography}{9}

\bibitem{Burkhart2026}
M.~C. Burkhart, \emph{Odd-ordered Cayley graphs satisfy Erd\H{o}s--Gy\'arf\'as}, arXiv:2611.05060 (2026).

\bibitem{Carr2026}
A.~Carr, \emph{Every minimal counterexample to the Erd\H{o}s--Gy\'arf\'as conjecture is predominantly cubic}, arXiv:2605.22844 (2026).

\bibitem{DucoffeDumitru2026}
G.~Ducoffe and B.~Dumitru, \emph{Towards a more structured search for Erd\H{o}s--Gy\'arf\'as counter-examples}, arXiv:2609.28594 (2026).

\bibitem{Erdos1997}
P.~Erd\H{o}s, \emph{Some old and new problems in various branches of combinatorics}, Discrete Math. \textbf{165/166} (1997), 227--231.

\bibitem{Garcia2026}
D.~Garcia, \emph{Small graphs without power-of-two cycles: a lower bound of 24, a correction to a construction of Exoo, and explicit bounds for $f(k)$}, arXiv:2609.04686 (2026).

\bibitem{LiuMontgomery2023}
H.~Liu and R.~Montgomery, \emph{A solution to Erd\H{o}s and Hajnal's odd cycle problem}, J. Amer. Math. Soc. \textbf{36} (2023), 1191--1234.

\bibitem{Tranquilli2026}
J.~Tranquilli, \emph{A 60-vertex lower bound for cubic bipartite counterexamples to the Erd\H{o}s--Gy\'arf\'as conjecture}, arXiv:2608.02675 (2026).

\end{thebibliography}

\end{document}
